\documentclass{article}
\usepackage[left=3cm, right=3cm, top=2cm]{geometry}

\usepackage{booktabs}
\usepackage[utf8]{inputenc}
\usepackage{tcolorbox}
\usepackage{enumitem}
\usepackage[obeyspaces]{url}
\usepackage{tabularx,array}
\newcolumntype{Y}{>{\raggedright\arraybackslash}X}
\newcolumntype{C}[1]{>{\centering\arraybackslash}p{#1}}

\usepackage{hyperref}
\hypersetup{
    colorlinks=true,
    linkcolor=blue,
    filecolor=magenta,
    urlcolor=cyan,
}
\usepackage{graphicx}
\usepackage{subcaption}
\usepackage{tikz}
\usetikzlibrary{automata, positioning, arrows}

\title{Artificial Intelligence Laboratory 2: A* Search Algorithm \\ DT8042 HT26, Halmstad University}
\author{Firstname Lastname, \{firstname.lastname\}@hh.se, \{programme\}}
\date{November 2026}

\begin{document}
\maketitle

\section*{Introduction}
The objective of this lab is to ...

\begin{tcolorbox}
\textbf{Q1}: What search algorithms have you learned from the lectures?
Briefly introduce them here.
\end{tcolorbox}

\section*{Task 1: Path Planning}
For Task 1 ...

\begin{tcolorbox}
\textbf{Q2}: Present your results for Random Search, BFS, and DFS on both map
types. Do they find optimal paths? How many nodes are expanded? Include
representative plots of the paths and the diagnostic values requested in the
lab description.
\end{tcolorbox}

\begin{tcolorbox}
\textbf{Q3}: Present your results for Greedy Search and A*. Which heuristics
have you implemented, including the custom heuristic for the rotated-H map?
Do the methods find optimal paths? Explain your observations with examples.
\end{tcolorbox}

\begin{tcolorbox}
\textbf{Q4}: Compare the algorithms over repeated trials (for example, 20)
and complete Table~\ref{tab:path-comparison}. State the map settings and number
of trials. Present results for the two map types separately, and add rows for
the heuristic variants already tested in the lab.
\end{tcolorbox}

\begin{table}[htbp]
\centering
\small
\setlength{\tabcolsep}{4pt}
\renewcommand{\arraystretch}{1.2}
\begin{tabularx}{\linewidth}{Y C{1.9cm} C{1.8cm} C{2.1cm} C{2cm}}
\toprule
Search algorithm & Nodes expanded & Path length (moves) & Time (s) & Optimal runs \\
\midrule
Random & \ldots & \ldots & \ldots & \ldots \\
BFS & \ldots & \ldots & \ldots & \ldots \\
DFS & \ldots & \ldots & \ldots & \ldots \\
Greedy & \ldots & \ldots & \ldots & \ldots \\
A* & \ldots & \ldots & \ldots & \ldots \\
\ldots & \ldots & \ldots & \ldots & \ldots \\
\bottomrule
\end{tabularx}
\caption{Performance comparison (Path Planning): [map type and settings].}
\label{tab:path-comparison}
\end{table}

Report repeated measurements as mean $\pm$ standard deviation. Path length
is the number of moves; average it over successful runs only. ``Optimal runs''
is the number of runs returning a shortest path, out of the solvable trials.
State failures separately. Time refers to the search, excluding plotting and
map generation.

\section*{Task 2: Poker Bidding}

\begin{tcolorbox}
\textbf{Q5}: Implement and apply BFS and DFS (and, optionally, Greedy Search).
What have you observed? Did each algorithm find a successful solution? How many
nodes were expanded?

\medskip
\textit{Task 2a:} What is the total number of states before the
fifth hand?
\end{tcolorbox}

\begin{tcolorbox}
\textbf{Q6} \textit{(Task 2c)}: Briefly describe your evaluation
function. How many nodes were expanded using it? Is the solution optimal?
If discussing optimality, state the criterion used.
\end{tcolorbox}

\begin{tcolorbox}
\textbf{Q7}: Compare BFS and DFS using Table~\ref{tab:poker-comparison} and
discuss the results. Include heuristic-search methods only if you completed
Task 2c.
\end{tcolorbox}

\begin{table}[htbp]
\centering
\small
\setlength{\tabcolsep}{4pt}
\renewcommand{\arraystretch}{1.2}
\begin{tabularx}{\linewidth}{Y C{1.5cm} C{2.3cm} C{2.3cm} C{2.1cm}}
\toprule
Search algorithm & Success? & Nodes expanded & Completed hands & Total bets \\
\midrule
BFS & \ldots & \ldots & \ldots & \ldots \\
DFS & \ldots & \ldots & \ldots & \ldots \\
\ldots & \ldots & \ldots & \ldots & \ldots \\
\bottomrule
\end{tabularx}
\caption{Performance comparison (Poker Bidding).}
\label{tab:poker-comparison}
\end{table}

Success is Yes/No according to the lab's winning objective.
Expanded nodes refer to the search. Hands and bets refer to the returned
successful sequence; use a dash if none is found. Total bets counts betting
actions by both players across all hands in that sequence, excluding calls
and folds. The supplied bidding counter resets after each hand.

\section*{Contribution}
Explain who worked on which parts.

\section*{Conclusion}
Add your reflections and conclusions about the lab.

\end{document}
